\section{Resumen ejecutivo}

Este documento establece la línea base conceptual para una plataforma interna de apoyo a la prevención cardiovascular. La propuesta busca transformar datos disponibles en conocimiento sobre factores asociados a enfermedad cardíaca y, en una etapa posterior, apoyar la priorización de evaluaciones preventivas. No corresponde a un sistema de diagnóstico autónomo ni reemplaza la decisión de los profesionales de salud.

Se comparan tres alternativas técnicamente equivalentes: una implementación Cloud en Google Cloud Platform (GCP), una alternativa Local dentro de la infraestructura hospitalaria y una arquitectura Híbrida. Las tres conservan el mismo alcance funcional, el mismo horizonte de evaluación y los mismos supuestos operacionales; cambian principalmente la ubicación de los componentes, la responsabilidad de operación, el perfil de costos y los riesgos.

La línea base define un stack portable y predominantemente open source, formado por Python, Pandas, Scikit-learn, PostgreSQL, SHAP, FastAPI, Streamlit, Docker y Git. La gestión propuesta es adaptativa, con ciclos incrementales, revisiones, hitos y capacidad de ajuste ante la incertidumbre propia de los datos, el modelamiento, los riesgos y la validación.

El piloto se evalúa a 36 meses, con 20 usuarios autorizados en el primer año, crecimiento a 25 y 30 usuarios, y concurrencia de 5, 6 y 8 usuarios respectivamente. Se asume una carga histórica inicial de 10.000 registros analíticos, 5.000 registros nuevos o actualizados durante el primer año y crecimiento anual de 10\%. La carga se actualiza mediante un proceso batch diario y la inferencia se solicita bajo demanda. La plataforma debe mantener una disponibilidad objetivo de 99\% durante el horario operacional, con RTO máximo de 8 horas, RPO máximo de 24 horas, respaldos diarios y pruebas mensuales de restauración.

La alternativa final no se declara en esta etapa. La recomendación deberá integrar la evaluación financiera, el análisis de riesgos, las restricciones éticas y legales, y la factibilidad técnica.
